\section{误差累积与在线修正}

\subsection{Compounding errors 的根源}

模仿学习的核心 challenge 在于，预测的动作 $\hat{a}$ 会改变下一个状态 $s'$，从而带来非 i.i.d. 的输入分布。小错会逐步被放大，最终产生 catastrophic failure。

\begin{warningbox}{分布偏移与 covariate shift}
专家状态分布 $p_{\text{expert}}(s)$ 与学习策略分布 $p_{\pi}(s)$ 不一致是系统性风险，如果不干预，即便均值训练 loss 很低也会在少数状态扯断 performance。
\end{warningbox}

\subsection{DAgger 与 Human-Gated DAgger}

解决思路是让策略自己运行，收集新的状态，并用专家在这些状态上提供修正动作：
\begin{enumerate}
    \item Rollout 当前策略 $\pi_\theta$，记录访问的状态。
    \item 在每个状态上请求专家动作，并补充到 dataset。
    \item 重新训练策略或通过 importance weighting 更新参数。
\end{enumerate}

\begin{table}[H]
\centering
\begin{tabular}{p{0.28\textwidth} p{0.28\textwidth} p{0.28\textwidth}}
\toprule
阶段 & 关键交付物 & 审查项 \\
\midrule
Observe & Rollout log & state distribution shift \\
Assess & Expert correction & human agreement rate \\
Act & Updated policy checkpoint & performance vs baseline \\
\bottomrule
\end{tabular}
\caption{DAgger 流程中的交付与审查}
\end{table}

\begin{knowledgebox}{Offline vs. Online 模仿学习}
\begin{itemize}
    \item \textbf{Offline（行为克隆）}：只使用 static dataset，数据安全但泛化差。
    \item \textbf{Online（DAgger/HG-DAgger）}：涉及 human-in-the-loop，效率高但需要专家权限。
\end{itemize}
\end{knowledgebox}

\subsection{模仿学习与强化学习的协同}

实践中常常把 BC 作为 warm-start，再用 RL fine-tune，形成 hybrid pipeline：先用高质量 demonstrations 学习初始 policy，再通过分布式 rollout + reward fine-tuning 提升鲁棒性。

\begin{warningbox}{不要放弃 reward 信号}
即便目标是模仿专家，也应持续监测潜在 reward（例如 safety penalty、human critique score），以便在奖励出现 drift 时立刻 rollback。
\end{warningbox}

\subsection{本章小结}

Compounding errors 需要数据收集与 human-in-the-loop 策略的结合。DAgger/HG-DAgger 提供了修正路径，结合后续 RL 微调可以让 policy 在 distribution shift 中继续收敛。

